\clearpage\pagestyle{otMain}\setlength{\otBlocEspacement}{2.5mm}\fontsize{11}{15}\selectfont\corrigesChapitre{6}{Les fractions, du tout au nombre}
\begin{blocLivre}{cor-6-1}\begin{corrige}{1}{Lire des parts}\demarche{Prendre le soixantième comme part.}
Une minute est $\frac1{60}$ d’heure : 7 minutes font $\res{\frac7{60}}$ d’heure et 45 minutes $\res{\frac{45}{60}}$, soit trois quarts d’heure.
\end{corrige}\end{blocLivre}
\begin{blocLivre}{cor-6-2}\begin{corrige}{2}{Représenter}\demarche{Compter les parts du partage.}
On colorie $\res5$ parts sur les 12 ; la solution figure en vert. Tout choix de cinq parts convient.
\label{fig-cor-6-2}\figureGeo[corrige]{frac-2}
\end{corrige}\end{blocLivre}
\begin{blocLivre}{cor-6-3}\begin{corrige}{3}{Entier caché}\demarche{Chercher les paquets complets du dénominateur.}
$18$ sixièmes font 3 unités : $\frac{19}6=\res{3+\frac16}$. Et $\frac{48}8=\res6$ exactement.
\end{corrige}\end{blocLivre}
\begin{blocLivre}{cor-6-4}\begin{corrige}{4}{Recomposer}\demarche{Convertir les entiers dans la part utilisée.}
$4=\frac{40}{10}$ donc $4+\frac3{10}=\res{\frac{43}{10}}$. De même $2=\frac{12}6$ et $2+\frac56=\res{\frac{17}6}$.
\end{corrige}\end{blocLivre}
\begin{blocLivre}{cor-6-5}\begin{corrige}{5}{Encadrer}\demarche{Encadrer le numérateur par des multiples du dénominateur.}
$\res{4<\frac{47}{10}<5}$, car $40<47<50$. Et $\res{5<\frac{23}4<6}$, car $20<23<24$.
\end{corrige}\end{blocLivre}
\begin{blocLivre}{cor-6-6}\begin{corrige}{6}{Lire la demi-droite}\demarche{Compter les unités puis les pas.}
A est 2 pas après 1 : $\res{1+\frac26=\frac86}$. B est 1 pas après 3 : $\res{3+\frac16=\frac{19}6}$.
\label{fig-cor-6-6}\figureGeo[corrige]{fd-6}
\end{corrige}\end{blocLivre}
\begin{blocLivre}{cor-6-7}\begin{corrige}{7}{Placer des fractions}\demarche{Convertir en entier plus fraction avant de placer.}
C se place 3 pas après 1, D trois pas après 2. Les positions figurent en vert au corrigé.
\label{fig-cor-6-7}\figureGeo[corrige]{fd-7}
\end{corrige}\end{blocLivre}
\begin{blocLivre}{cor-6-8}\begin{corrige}{8}{Le même point}\demarche{Compter les pas couverts par chaque écriture.}
$\res{\text{Oui}}$ : 6 pas de un douzième couvrent la moitié de l’unité, exactement comme $\frac12$. Deux écritures différentes peuvent désigner le même nombre.
\label{fig-cor-6-8}\figureGeo[corrige]{fd-8}
\end{corrige}\end{blocLivre}
\begin{blocLivre}{cor-6-9}\begin{corrige}{9}{Un entier tout rond}\demarche{Multiplier l’entier par le dénominateur.}
$2=\frac{30}{15}$ : $\square=\res{30}$. $36:9=4$ : $\square=\res9$. Une unité entière : $\square=\res{60}$.
\end{corrige}\end{blocLivre}
\begin{blocLivre}{cor-6-10}\begin{corrige}{10}{Les tours de piste}\demarche{Extraire les tours entiers.}
$\frac{11}4=2+\frac34$ : $\res2$ tours complets et trois quarts du troisième. Il lui reste $\res{\frac14}$ de tour pour finir le troisième.
\end{corrige}\end{blocLivre}
\begin{blocLivre}{cor-6-11}\begin{corrige}{11}{Les gobelets}\demarche{Compter directement les quarts.}
Treize quarts remplissent $\res{13}$ gobelets d’un quart. Après 3 gobelets, il reste $\frac{13}4-\frac34=\frac{10}4$ de litre, soit $\res{2+\frac24}$ litres.
\end{corrige}\end{blocLivre}
\begin{blocLivre}{cor-6-12}\begin{corrige}{12}{Entre deux entiers}\demarche{Convertir les bornes en sixièmes.}
Entre $\frac{12}6$ et $\frac{18}6$, on peut choisir $\res{\frac{13}6,\ \frac{14}6,\ \frac{17}6}$ par exemple : tout numérateur de 13 à 17 convient.
\end{corrige}\end{blocLivre}
\begin{blocLivre}{cor-6-13}\begin{corrige}{13}{Le pot de peinture}\demarche{Lire l’entier puis interpréter le reste.}
$\frac{35}{10}=\res{3+\frac5{10}}$ : trois pots entiers et la moitié d’un quatrième. Il a ouvert au minimum $\res4$ pots.
\end{corrige}\end{blocLivre}
\begin{blocLivre}{cor-6-14}\begin{corrige}{14}{Comparer à l’unité et à 2}\demarche{Comparer le numérateur à 9 et à 18.}
$\frac79<1$ ; $\frac99=1$ ; $\frac{13}9$ est entre 1 et 2, car $9<13<18$ ; $\frac{25}9>2$, car $25>18$. Groupes : $\res{\frac79}$ ; $\res{\frac{13}9}$ ; $\res{\frac{25}9}$ ; et $\frac99$ vaut exactement 1.
\end{corrige}\end{blocLivre}
\begin{blocLivre}{cor-6-15}\begin{corrige}{15}{Vrai ou faux ?}\demarche{Comparer chaque fraction à l’unité.}
Les rôles diffèrent : $\frac{60}{59}$ dépasse 1 (60 parts d’un partage en 59), tandis que $\frac{59}{60}<1$ (il manque un soixantième). Elles sont $\res{\text{différentes}}$ : l’ordre des nombres compte.
\end{corrige}\end{blocLivre}
